\centerline{\bf \Huge Китайская теорема об остатках}


\begin{flushright}
        \scriptsize{Эпизод 20. Очертания сердца, очертания человека }
\end{flushright}


\problem{1}
Генерал построил солдат в колонну по 4 , но при этом солдат Иванов остался лишним. Тогда генерал построил солдат в колонну по 5 . И снова Иванов остался лишним. Когда же и в колонне по 6 Иванов оказался лишним, генерал посулил ему наряд вне очереди, после чего в колонне по 7 Иванов нашел себе место и никого лишнего не осталось. Какое наименьшее число солдат могло быть у генерала?

\problem{2}
Сколько четырёхзначных чисел подходит под условие $x^2 \equiv x(\bmod 10000)$ ?

\problem{3}
Диме выдали натуральное число $N$. Он разделил его на 101 и получил в остатке $m>0$. Затем Дима разделил $N$ на $m$ и получил в остатке $p$. Найдите наибольшее значение $p$, которое могло получиться, а затем - наименьшее $N$, при котором это значение $p$ достигается.

\problem{4}
Докажите, что натуральные $n$, для которых $n^n+1$ делится на 30 , образуют арифметическую прогрессию.

\problem{5}
Найдите все натуральные $n>1$, для которых любое целое число можно представить в виде суммы двух целых чисел, каждое из которых взаимно просто с $n$.

\problem{6}
Даны натуральные $a$ и $b$. Известно, что для любого натурального $n$ число $a^n+n$ делится на $b^n+n$. Докажите, что $a=b$.